\section{Insertion}
\label{sec:insertion}

Insertion into an AVL tree begins exactly as insertion into an ordinary binary
search tree. The new key descends from the root, comparing at each node, until it
reaches an empty position, where it becomes a leaf. What follows is the part
\Cref{def:avl} forces. The new leaf may have lengthened the path beneath some of
its ancestors, so the procedure walks back up, updating heights and testing
balance factors, and repairs the first node it finds in violation.

\begin{algorithm}[t]
\caption{Insertion. The recursion descends to the insertion point and rebalances
on the way back out, so every node on the search path is repaired after its
children and before its parent.}
\label{alg:insert}
\SetKwFunction{Ins}{Insert}
\SetKwFunction{Reb}{Rebalance}
\Fn{\Ins{$v$, $k$}}{
  \lIf{$v = \textbf{nil}$}{\KwRet \textbf{new} node with key $k$ and height $0$}
  \uIf{$k < v.\mathit{key}$}{
    $v.\mathit{left} \gets$ \Ins{$v.\mathit{left}$, $k$}\;
  }
  \uElseIf{$k > v.\mathit{key}$}{
    $v.\mathit{right} \gets$ \Ins{$v.\mathit{right}$, $k$}\;
  }
  \lElse{\KwRet $v$ \Comment*[f]{the key is already present}}
  \KwRet \Reb{$v$} \Comment*[r]{updates the height, rotates if needed}
}
\end{algorithm}

\Cref{fig:inserttrace} follows one insertion through the three stages. The key
$5$ arrives as a leaf under the node $10$. Walking back up, the nodes $10$ and
$20$ both reach a balance factor of $+1$, which \Cref{def:avl} permits, and the
node $40$ reaches $+2$, which it does not. The path from $40$ down toward the new
leaf goes left and then left again, so \Cref{tab:cases} calls for a single right
rotation at $40$. After that rotation the tree is legal, and the walk stops.

\begin{figure}[t]
  \centering
  % One insertion, start to finish. Key 5 arrives as a leaf, the balance factors
% are recomputed on the way back to the root, node 40 is the first to reach
% two, and the single rotation there settles the whole tree.
%
% Geometry is tight on purpose: three panels have to fit inside 158mm of text
% block. Panel pitch is 52mm and no label sticks out to the left of a panel.
\begin{tikzpicture}[x = 1mm, y = 1mm]

  % ------------------------------------------------- (a) after the plain insert
  \begin{scope}
    \node[tnode] (a40) at (  5,   0) {40};
    \node[tnode] (a20) at ( -6, -12) {20};
    \node[tnode] (a60) at ( 16, -12) {60};
    \node[tnode] (a10) at (-14, -24) {10};
    \node[tnode] (a30) at (  2, -24) {30};
    \node[tmark] (a05) at (-21, -36) {5};
    \foreach \p/\c in {a40/a20, a40/a60, a20/a10, a20/a30, a10/a05}
      \draw[tedge] (\p) -- (\c);
    \node[tlabel, below = 1mm of a05] {new};
    \node[tannot] at (0, -48) {(a) inserted as a leaf};
  \end{scope}

  % ------------------------------------------ (b) balance factors on the way up
  \begin{scope}[xshift = 52mm]
    \node[tmark]  (b40) at (  5,   0) {40};
    \node[tnode]  (b20) at ( -6, -12) {20};
    \node[tnode]  (b60) at ( 16, -12) {60};
    \node[tnode]  (b10) at (-14, -24) {10};
    \node[tnode]  (b30) at (  2, -24) {30};
    \node[tnode]  (b05) at (-21, -36) {5};
    \foreach \p/\c in {b40/b20, b40/b60, b20/b10, b20/b30, b10/b05}
      \draw[tedge] (\p) -- (\c);
    \draw[tedge, line width = 1.2pt] (b40) -- (b20);
    \draw[tedge, line width = 1.2pt] (b20) -- (b10);
    \draw[tedge, line width = 1.2pt] (b10) -- (b05);
    % all three labels hang to the right, so nothing overhangs the panel
    \node[tlabel, right = 1mm of b40] {$+2$};
    \node[tlabel, right = 1mm of b20] {$+1$};
    \node[tlabel, right = 1mm of b10] {$+1$};
    \node[tannot] at (0, -48) {(b) node 40 reaches $+2$};
  \end{scope}

  % ------------------------------------------------------ (c) after the rotation
  \begin{scope}[xshift = 104mm]
    \node[tnode] (c20) at (  0,  -4) {20};
    \node[tnode] (c10) at (-12, -17) {10};
    \node[tnode] (c40) at ( 12, -17) {40};
    \node[tnode] (c05) at (-20, -30) {5};
    \node[tnode] (c30) at (  3, -30) {30};
    \node[tnode] (c60) at ( 20, -30) {60};
    \foreach \p/\c in {c20/c10, c20/c40, c10/c05, c40/c30, c40/c60}
      \draw[tedge] (\p) -- (\c);
    \node[tlabel, right = 1mm of c20] {$0$};
    \node[tannot] at (0, -48) {(c) one right rotation at 40};
  \end{scope}

\end{tikzpicture}

  \caption{One insertion. The key $5$ enters as a leaf, the node $40$ is the
  lowest ancestor to reach a balance factor of $+2$, and one right rotation there
  restores the invariant for the whole tree.}
  \label{fig:inserttrace}
\end{figure}

The walk stopped after one repair, and that is not particular to this example.

\begin{theorem}[One repair suffices]
\label{thm:insertone}
An insertion into an AVL tree triggers at most one rebalancing step, consisting
of either one single rotation or one double rotation.
\end{theorem}

\begin{proof}
Let $z$ be the lowest ancestor of the new leaf whose balance factor reaches
$\pm 2$, and suppose the subtree rooted at $z$ had height $h$ before the
insertion. An insertion adds one leaf, so it increases the height of any subtree
by at most one, and $z$ became unbalanced only because the height of its taller
child grew. The subtree at $z$ therefore has height $h+1$ after the insertion and
before the repair.

Apply the repair \Cref{tab:cases} prescribes for the case at hand. In the
\case{ll} case, write $y$ for the left child of $z$ and let $A = L(y)$, $B =
R(y)$ and $C = R(z)$ as in the proof of \Cref{lem:inorder}. Since $z$ has balance
factor $+2$ and $y$ has balance factor $+1$, the subtree $A$ has height $h-1$
while $B$ and $C$ both have height $h-2$. After the right rotation at $z$, the
node $y$ roots the subtree with $A$ on its left and, on its right, the node $z$
carrying $B$ and $C$. That right side has height $h-1$, matching $A$, so the
rotated subtree has height $h$. The \case{rr} case is the mirror image, and the
two double-rotation cases give the same result by the same count.

The subtree rooted at the repaired position therefore has height $h$ again, which
is the height it had before the insertion. Every proper ancestor of $z$ sees a
child of unchanged height, so no ancestor changes its own height and no ancestor
changes its balance factor. No further node can be in violation, and the walk
terminates.
\end{proof}

\Cref{fig:whystops} states the argument in one picture. The insertion adds a
level, the rotation takes that level back, and nothing above the repaired node
observes either event.

\begin{corollary}[Insertion cost]
\label{cor:insertcost}
Insertion into an AVL tree of $n$ nodes takes $O(\log n)$ time.
\end{corollary}

\begin{proof}
The descent visits one node per level, and by \Cref{thm:heightbound} there are
$O(\log n)$ levels. The upward walk visits the same nodes, spending constant time
on each to update a height and test a balance factor. By \Cref{thm:insertone} at
most one rebalancing step runs, and by \Cref{lem:rotcost} it costs $O(1)$.
Summing the three parts gives $O(\log n)$.
\end{proof}

\begin{figure}[t]
  \centering
  \begin{tikzpicture}[x = 1mm, y = 1mm]
    \foreach \dx/\hh/\lab in {0/22/{before}, 46/28/{after the insertion},
                              92/22/{after the repair}} {
      \begin{scope}[xshift = \dx mm]
        \node[tnode, minimum size = 4mm] (r\dx) at (0, 0) {};
        \draw[tedge] (r\dx) -- (-12, -\hh);
        \draw[tedge] (r\dx) -- ( 12, -\hh);
        \draw[tedge] (-12, -\hh) -- (12, -\hh);
        \node[tannot] at (0, -\hh - 7) {\lab};
      \end{scope}
    }
    \draw[-{Stealth[length=2.5mm]}, draw = rmid] (16, -14) -- (30, -14)
      node[midway, above, tlabel] {insert};
    \draw[-{Stealth[length=2.5mm]}, draw = rmid] (62, -14) -- (76, -14)
      node[midway, above, tlabel] {rotate};
    \draw[decorate, decoration = {brace, amplitude = 3pt}, draw = rmid]
      (-16, -22) -- (-16, 0)
      node[midway, left = 4pt, tlabel] {$h$};
    \draw[decorate, decoration = {brace, amplitude = 3pt, mirror}, draw = rmid]
      (108, -22) -- (108, 0)
      node[midway, right = 4pt, tlabel] {$h$};
  \end{tikzpicture}
  \caption{Why one repair ends the walk. The insertion raises the subtree by one
  level and the rotation returns it, so the parent of the repaired node sees no
  change and the invariant above it was never disturbed.}
  \label{fig:whystops}
\end{figure}
